\subjudul{Kontrak, class pembawa data, dan cara Python menilai sebuah tipe.}
\chapter{Abstract Method, Data Class, dan Typing}

\begin{dibabini}
  \item Abstract method
  \item Data class
  \item Attribute mutability
  \item Duck typing
  \item Structural typing
\end{dibabini}

Di Praktik Mandiri Bab 4, method \texttt{luas()} pada class \texttt{BangunDatar} mengembalikan \texttt{0}. Cara itu menyimpan masalah: jika seseorang membuat class \texttt{Segitiga} lalu lupa meng-override \texttt{luas()}, program tidak error, tetapi diam-diam menghasilkan luas nol. Kita butuh cara untuk \emph{mewajibkan} child menyediakan method tertentu.

Bab ini membahas lima topik yang saling berkaitan. \emph{Abstract method} dipakai untuk membuat kontrak yang wajib dipenuhi child class. \emph{Data class} menghemat penulisan class yang tugas utamanya membawa data. \emph{Mutability} menjelaskan mengapa sebagian data bisa berubah ``dari jauh''. Terakhir, \emph{duck typing} dan \emph{structural typing} menjelaskan cara Python menilai apakah sebuah object cocok dipakai di suatu tempat.

% =====================================================================
\section{Abstract Method}
% =====================================================================

\emph{Abstract method} adalah method yang hanya dideklarasikan namanya di parent, tanpa implementasi yang berarti, dan \textbf{wajib} diimplementasikan oleh setiap child class. Class yang memiliki setidaknya satu abstract method disebut \emph{abstract class}. Abstract class tidak bisa dibuat object-nya secara langsung; ia hanya berfungsi sebagai \emph{kontrak} atau cetakan bagi turunannya.

\subsection{Masalah: Lupa Meng-override}

\begin{kode}[Parent dengan nilai default yang menyesatkan]
class BangunDatar:
    def luas(self):
        return 0                      # nilai "darurat"

class Segitiga(BangunDatar):
    def __init__(self, alas, tinggi):
        self.alas = alas
        self.tinggi = tinggi
    # lupa menulis luas()

print(Segitiga(10, 6).luas())
\end{kode}

\begin{keluaran}
0
\end{keluaran}

Tidak ada pesan kesalahan sama sekali, padahal hasilnya jelas salah. Kesalahan seperti ini bisa lolos hingga ke tangan pengguna.

\subsection{Modul \texttt{abc}: \texttt{ABC} dan \texttt{@abstractmethod}}

Python menyediakan modul \texttt{abc} (\emph{Abstract Base Classes}). Sebuah class menjadi abstract class bila mewarisi \texttt{ABC}, dan method-nya ditandai sebagai abstract dengan decorator \texttt{@abstractmethod}.

\begin{kode}[Abstract class BangunDatar]
from abc import ABC, abstractmethod
import math

class BangunDatar(ABC):
    def __init__(self, nama):
        self.nama = nama

    @abstractmethod
    def luas(self):
        """Wajib diimplementasikan child."""

    @abstractmethod
    def keliling(self):
        """Wajib diimplementasikan child."""

    def info(self):                       # method biasa (konkret)
        return f"{self.nama}: luas {self.luas():.2f}, keliling {self.keliling():.2f}"

class Persegi(BangunDatar):
    def __init__(self, sisi):
        super().__init__("Persegi")
        self.sisi = sisi

    def luas(self):
        return self.sisi ** 2

    def keliling(self):
        return 4 * self.sisi

class Lingkaran(BangunDatar):
    def __init__(self, r):
        super().__init__("Lingkaran")
        self.r = r

    def luas(self):
        return math.pi * self.r ** 2

    def keliling(self):
        return 2 * math.pi * self.r

for b in [Persegi(5), Lingkaran(7)]:
    print(b.info())
\end{kode}

\begin{keluaran}
Persegi: luas 25.00, keliling 20.00
Lingkaran: luas 153.94, keliling 43.98
\end{keluaran}

Kini Python menjaga kontrak tersebut. Abstract class tidak bisa dibuat object-nya, dan child yang belum mengimplementasikan \emph{semua} abstract method juga ditolak.

\begin{kode}[Kontrak yang dilanggar]
class Segitiga(BangunDatar):
    def __init__(self, alas, tinggi):
        super().__init__("Segitiga")
        self.alas = alas
        self.tinggi = tinggi

    def luas(self):
        return 0.5 * self.alas * self.tinggi
    # keliling() belum ditulis

BangunDatar("Bebas")     # gagal: abstract class
Segitiga(10, 6)          # gagal: kontrak belum lengkap
\end{kode}

\begin{keluaran}
TypeError: Can't instantiate abstract class BangunDatar without an implementation for abstract methods 'keliling', 'luas'
TypeError: Can't instantiate abstract class Segitiga without an implementation for abstract method 'keliling'
\end{keluaran}

Setiap baris di atas sebenarnya menghentikan program; keduanya ditampilkan bersamaan agar mudah dibandingkan. Yang penting, kesalahan kini muncul \emph{saat object dibuat}, jauh sebelum nilai yang salah sempat dipakai.

\begin{figure}[htbp]
\centering
\begin{tikzpicture}[font=\sffamily\small,
    abs/.style={rectangle,draw=marun,line width=1pt,dashed,fill=merahmuda!50,align=left,text width=46mm,inner sep=2.5mm},
    kon/.style={rectangle,fill=marun,text=white,align=left,text width=30mm,inner sep=2.5mm},
    gagal/.style={rectangle,draw=jingga,line width=1pt,fill=jingga!12,align=left,text width=30mm,inner sep=2.5mm}]
  \node[abs] (p) at (0,0)
    {\textit{\textbf{BangunDatar}} {\footnotesize\textcolor{abu}{(abstract)}}\\[3pt]
     \ttfamily\footnotesize \textit{+ luas()}\\\textit{+ keliling()}\\+ info()};
  \node[kon,anchor=north] (a) at (-4.4,-2.6) {\textbf{Persegi}\\[2pt]\ttfamily\footnotesize + luas()\\+ keliling()};
  \node[kon,anchor=north] (b) at (0,-2.6)    {\textbf{Lingkaran}\\[2pt]\ttfamily\footnotesize + luas()\\+ keliling()};
  \node[gagal,anchor=north] (c) at (4.4,-2.6) {\textbf{Segitiga}\\[2pt]\ttfamily\footnotesize + luas()\\\textcolor{jingga}{keliling()?}};
  \coordinate (cab) at (0,-1.9);
  \draw[marun,line width=0.9pt] (a.north) |- (cab) -| (c.north);
  \draw[marun,line width=0.9pt] (b.north) -- (cab);
  \draw[turunan,marun] (cab) -- (p.south);
  \node[anchor=north,text=hijau,font=\sffamily\footnotesize] at (a.south) {kontrak lengkap};
  \node[anchor=north,text=hijau,font=\sffamily\footnotesize] at (b.south) {kontrak lengkap};
  \node[anchor=north,text=jingga,font=\sffamily\footnotesize,align=center] at (c.south) {belum lengkap:\\object tidak bisa dibuat};
\end{tikzpicture}
\caption{Abstract class sebagai kontrak. Method bercetak miring wajib diimplementasikan setiap child; \texttt{info()} adalah method konkret yang langsung diwarisi.}
\end{figure}

Perhatikan bahwa abstract class boleh memiliki method biasa seperti \texttt{info()} dan constructor. Method konkret tersebut bahkan boleh memanggil abstract method (\texttt{self.luas()}), karena dijamin setiap object yang berhasil dibuat sudah mengimplementasikannya.

\subsection{Abstract Property}

Decorator \texttt{@property} dan \texttt{@abstractmethod} dapat digabung untuk mewajibkan child menyediakan sebuah property. Urutannya: \texttt{@property} di atas, \texttt{@abstractmethod} di bawah.

\begin{kode}[Abstract property]
from abc import ABC, abstractmethod

class Kendaraan(ABC):
    @property
    @abstractmethod
    def jumlah_roda(self):
        ...

class Motor(Kendaraan):
    @property
    def jumlah_roda(self):
        return 2

print(Motor().jumlah_roda)
\end{kode}

\begin{keluaran}
2
\end{keluaran}

\begin{info}[Abstract Class dan Interface]
Bahasa seperti Java memiliki konsep \emph{interface}, yaitu kontrak berisi daftar method tanpa implementasi. Di Python, abstract class yang seluruh method-nya abstract memainkan peran yang sama. Badan abstract method cukup diisi docstring atau \texttt{...} (\emph{ellipsis}).
\end{info}

% =====================================================================
\section{Data Class}
% =====================================================================

Banyak class yang tugas utamanya hanya \emph{membawa data}: koordinat titik, data mahasiswa, item keranjang belanja, dan sebagainya. Untuk class seperti itu, kita hampir selalu menulis tiga method yang sama polanya: \texttt{\_\_init\_\_} untuk menyimpan data, \texttt{\_\_repr\_\_} agar object mudah dicetak, dan \texttt{\_\_eq\_\_} agar dua object berisi data sama dianggap sama.

\begin{kode}[Class pembawa data yang ditulis manual]
class Titik:
    def __init__(self, x, y):
        self.x = x
        self.y = y

    def __repr__(self):
        return f"Titik(x={self.x}, y={self.y})"

    def __eq__(self, lain):
        if not isinstance(lain, Titik):
            return NotImplemented
        return (self.x, self.y) == (lain.x, lain.y)
\end{kode}

Lima belas baris hanya untuk menyimpan dua angka. Modul \texttt{dataclasses} menyediakan decorator \texttt{@dataclass} yang menuliskan method-method tersebut secara otomatis berdasarkan daftar \emph{field} yang kita deklarasikan.

\subsection{Decorator \texttt{@dataclass}}

Field ditulis sebagai class attribute yang diberi \emph{type hint} (anotasi tipe, dibahas lebih lanjut di Subbab Structural Typing). Dari daftar field itulah \texttt{@dataclass} membangkitkan \texttt{\_\_init\_\_}, \texttt{\_\_repr\_\_}, dan \texttt{\_\_eq\_\_}.

\begin{kode}[Titik sebagai data class]
from dataclasses import dataclass

@dataclass
class Titik:
    x: float
    y: float

    def jarak_ke_pusat(self):              # method biasa tetap boleh
        return (self.x ** 2 + self.y ** 2) ** 0.5

a = Titik(3, 4)
b = Titik(3, 4)
print(a)
print(a == b)
print(a.jarak_ke_pusat())
\end{kode}

\begin{keluaran}
Titik(x=3, y=4)
True
5.0
\end{keluaran}

\begin{figure}[htbp]
\centering
\begin{tikzpicture}[font=\sffamily\small,
    gen/.style={rectangle,draw=marun,line width=0.8pt,fill=merahmuda,font=\ttfamily\small,
      text width=62mm,align=left,inner sep=2.2mm}]
  \node[rectangle,fill=kodebg,text=kodeteks,font=\ttfamily\small,align=left,inner sep=3mm] (src) at (0,0)
    {\textcolor{emas}{@dataclass}\\\textcolor{emas}{class} Titik:\\\quad x: float\\\quad y: float};
  \node[gen,anchor=west] (g1) at (4.3,1.3)  {def \_\_init\_\_(self, x, y)};
  \node[gen,anchor=west] (g2) at (4.3,0)    {def \_\_repr\_\_(self)};
  \node[gen,anchor=west] (g3) at (4.3,-1.3) {def \_\_eq\_\_(self, lain)};
  \foreach \g in {g1,g2,g3}
    \draw[-{Latex},line width=1pt,emas] (src.east) -- (\g.west);
  \node[anchor=south,font=\sffamily\footnotesize,text=abu] at (6.6,1.9) {dibuatkan otomatis};
\end{tikzpicture}
\caption{Dari daftar field, \texttt{@dataclass} membangkitkan method-method pembawa data secara otomatis.}
\end{figure}

\subsection{Nilai Default dan \texttt{field()}}

Field boleh diberi nilai default. Seperti parameter fungsi, field ber-default harus ditulis setelah field tanpa default. Untuk default berupa object \emph{mutable} seperti list atau dictionary, Python melarang penulisan langsung; gunakan \texttt{field(default\_factory=\ldots)} agar setiap object mendapat salinan baru. Alasannya dibahas di subbab Attribute Mutability.

\begin{kode}[Nilai default]
from dataclasses import dataclass, field

@dataclass
class Mahasiswa:
    nama: str
    nim: str
    semester: int = 1
    nilai: list = field(default_factory=list)   # list baru untuk tiap object

ani = Mahasiswa("Ani", "001")
budi = Mahasiswa("Budi", "002", semester=3)
ani.nilai.append(90)
print(ani)
print(budi)
\end{kode}

\begin{keluaran}
Mahasiswa(nama='Ani', nim='001', semester=1, nilai=[90])
Mahasiswa(nama='Budi', nim='002', semester=3, nilai=[])
\end{keluaran}

Jika kita menulis \texttt{nilai: list = []}, Python langsung menolak saat class didefinisikan dengan pesan \texttt{ValueError: mutable default <class 'list'> for field nilai is not allowed: use default\_factory}.

\subsection{Opsi \texttt{order} dan \texttt{frozen}}

Decorator \texttt{@dataclass} menerima beberapa opsi. Dua yang paling sering dipakai:

\begin{itemize}
  \item \texttt{order=True}: membangkitkan \texttt{<}, \texttt{<=}, \texttt{>}, \texttt{>=}, sehingga object bisa diurutkan. Perbandingan dilakukan field demi field sesuai urutan deklarasi.
  \item \texttt{frozen=True}: membuat object \emph{tidak bisa diubah} setelah dibuat. Setiap upaya mengubah field menghasilkan \texttt{FrozenInstanceError}.
\end{itemize}

\begin{kode}[order dan frozen]
from dataclasses import dataclass, field

@dataclass(order=True)
class Peserta:
    skor: int
    nama: str = field(compare=False)      # tidak ikut dibandingkan

@dataclass(frozen=True)
class Koordinat:
    lintang: float
    bujur: float

hasil = [Peserta(78, "Ani"), Peserta(92, "Budi"), Peserta(85, "Citra")]
for p in sorted(hasil, reverse=True):
    print(p.skor, p.nama)

unesa = Koordinat(-7.30, 112.67)
unesa.lintang = 0
\end{kode}

\begin{keluaran}
92 Budi
85 Citra
78 Ani
Traceback (most recent call last):
  ...
dataclasses.FrozenInstanceError: cannot assign to field 'lintang'
\end{keluaran}

\subsection{Validasi dengan \texttt{\_\_post\_init\_\_}}

Karena \texttt{\_\_init\_\_} dibuat otomatis, validasi dan field turunan ditulis di method khusus \texttt{\_\_post\_init\_\_}, yang dijalankan tepat setelah \texttt{\_\_init\_\_} selesai. Field yang nilainya dihitung (bukan diberikan pengguna) ditandai \texttt{field(init=False)}.

\begin{kode}[Validasi dan field turunan]
from dataclasses import dataclass, field

@dataclass
class ItemBelanja:
    nama: str
    harga: int
    jumlah: int = 1
    subtotal: int = field(init=False)

    def __post_init__(self):
        if self.harga <= 0 or self.jumlah <= 0:
            raise ValueError("Harga dan jumlah harus lebih dari nol.")
        self.subtotal = self.harga * self.jumlah

print(ItemBelanja("Pulpen", 3500, 4))
ItemBelanja("Buku", -1)
\end{kode}

\begin{keluaran}
ItemBelanja(nama='Pulpen', harga=3500, jumlah=4, subtotal=14000)
Traceback (most recent call last):
  ...
ValueError: Harga dan jumlah harus lebih dari nol.
\end{keluaran}

\begin{table}[htbp]
\caption{Ringkasan fitur data class.}
\begin{tabularx}{\linewidth}{P{58mm}L}
\kepala \kp{Penulisan} & \kp{Hasilnya} \\
\texttt{@dataclass} & \texttt{\_\_init\_\_}, \texttt{\_\_repr\_\_}, \texttt{\_\_eq\_\_} dibuat otomatis \\
\texttt{@dataclass(order=True)} & ditambah operator perbandingan sehingga bisa diurutkan \\
\texttt{@dataclass(frozen=True)} & object tidak bisa diubah setelah dibuat \\
\texttt{field(default\_factory=list)} & default mutable yang aman, baru untuk tiap object \\
\texttt{field(compare=False)} & field tidak ikut perbandingan \texttt{==} dan pengurutan \\
\texttt{field(init=False)} & field tidak menjadi parameter constructor \\
\texttt{\_\_post\_init\_\_} & tempat validasi dan menghitung field turunan \\
\end{tabularx}
\end{table}

\begin{tips}[Kapan Memakai Data Class?]
Gunakan data class untuk class yang \emph{intinya adalah data}, dengan sedikit perilaku. Untuk class yang kaya aturan dan perlu encapsulation ketat, misalnya \texttt{Rekening} dengan saldo yang hanya boleh berubah lewat \texttt{setor()} dan \texttt{tarik()}, class biasa dengan property sering kali lebih tepat.
\end{tips}

% =====================================================================
\section{Attribute Mutability}
% =====================================================================

\emph{Mutability} adalah sifat apakah isi sebuah object dapat diubah setelah object itu dibuat. Object \emph{mutable} dapat diubah isinya; object \emph{immutable} tidak dapat diubah, sehingga setiap ``perubahan'' sebenarnya menghasilkan object baru.

\begin{table}[htbp]
\caption{Tipe bawaan berdasarkan mutability.}
\begin{tabularx}{\linewidth}{P{30mm}L}
\kepala \kp{Sifat} & \kp{Tipe} \\
Immutable & \texttt{int}, \texttt{float}, \texttt{bool}, \texttt{str}, \texttt{tuple}, \texttt{frozenset}, data class dengan \texttt{frozen=True} \\
Mutable & \texttt{list}, \texttt{dict}, \texttt{set}, sebagian besar object dari class buatan sendiri \\
\end{tabularx}
\end{table}

\subsection{Variabel Adalah Label, Bukan Kotak}

Untuk memahami mutability, kita perlu meluruskan satu anggapan. Variabel di Python bukan ``kotak'' yang berisi nilai, melainkan \emph{label} yang ditempelkan pada sebuah object. Pernyataan \texttt{b = a} tidak menyalin object; ia hanya menempelkan label kedua pada object yang sama. Kondisi ini disebut \emph{aliasing}.

\begin{kode}[Dua label, satu object]
a = [1, 2, 3]
b = a              # b menunjuk object yang sama
b.append(4)
print(a)
print(a is b)

c = a.copy()       # object baru dengan isi sama
c.append(5)
print(a, c)
print(a is c)
\end{kode}

\begin{keluaran}
[1, 2, 3, 4]
True
[1, 2, 3, 4] [1, 2, 3, 4, 5]
False
\end{keluaran}

\begin{figure}[htbp]
\centering
\begin{tikzpicture}[font=\sffamily\small,
    label/.style={rectangle,fill=merah,text=white,font=\ttfamily\small\bfseries,minimum width=9mm,minimum height=7mm},
    sel/.style={rectangle,draw=marun,fill=merahmuda,font=\ttfamily\small,minimum width=8mm,minimum height=8mm,inner sep=0pt},
    tunjuk/.style={-{Latex},line width=1pt,emas}]
  % Panel kiri: b = a
  \node[font=\sffamily\small\bfseries,text=merah] at (1.6,2.4) {\texttt{b = a}};
  \node[label] (a1) at (0,1.2) {a};
  \node[label] (b1) at (0,-0.2) {b};
  \foreach \i/\v in {0/1,1/2,2/3,3/4} \node[sel] (s\i) at (2.0+0.8*\i,0.5) {\v};
  \draw[tunjuk] (a1.east) -- (s0.north west);
  \draw[tunjuk] (b1.east) -- (s0.south west);
  % Panel kanan: c = a.copy()
  \node[font=\sffamily\small\bfseries,text=merah] at (8.6,2.4) {\texttt{c = a.copy()}};
  \node[label] (a2) at (7,1.2) {a};
  \node[label] (c2) at (7,-0.6) {c};
  \foreach \i/\v in {0/1,1/2,2/3,3/4} \node[sel] (t\i) at (9.0+0.8*\i,1.2) {\v};
  \foreach \i/\v in {0/1,1/2,2/3,3/4,4/5} \node[sel] (u\i) at (9.0+0.8*\i,-0.6) {\v};
  \draw[tunjuk] (a2.east) -- (t0.west);
  \draw[tunjuk] (c2.east) -- (u0.west);
  \draw[garis,line width=0.6pt] (5.6,-1.2) -- (5.6,2.7);
\end{tikzpicture}
\caption{Kiri: \texttt{a} dan \texttt{b} adalah dua label pada satu list, sehingga perubahan lewat \texttt{b} terlihat lewat \texttt{a}. Kanan: \texttt{copy()} membuat list baru yang terpisah.}
\end{figure}

Untuk object immutable, aliasing tidak berbahaya. Operasi seperti \texttt{y += 1} pada sebuah \texttt{int} tidak mengubah object lama, melainkan membuat object baru dan memindahkan label \texttt{y} ke sana, sehingga variabel lain yang menunjuk object lama tidak terpengaruh.

\subsection{Jebakan: Default Argument Mutable}

Nilai default parameter fungsi dibuat \textbf{satu kali}, saat fungsi didefinisikan, bukan setiap kali fungsi dipanggil. Jika nilai default itu mutable, semua pemanggilan memakai object yang sama. Di dalam class, akibatnya semua object diam-diam berbagi satu list.

\begin{kode}[Bug default mutable]
class Keranjang:
    def __init__(self, pemilik, isi=[]):     # BAHAYA
        self.pemilik = pemilik
        self.isi = isi

    def tambah(self, barang):
        self.isi.append(barang)

k1 = Keranjang("Ani")
k2 = Keranjang("Budi")
k1.tambah("Buku")
print(k2.isi)            # keranjang Budi ikut berisi buku!
\end{kode}

\begin{keluaran}
['Buku']
\end{keluaran}

Perbaikan yang lazim adalah memakai \texttt{None} sebagai default, lalu membuat list baru di dalam fungsi.

\begin{kode}[Perbaikan]
class Keranjang:
    def __init__(self, pemilik, isi=None):
        self.pemilik = pemilik
        self.isi = [] if isi is None else list(isi)
\end{kode}

Perhatikan pemakaian \texttt{list(isi)}: walaupun pemanggil memberikan list miliknya sendiri, keranjang menyimpan \emph{salinannya}. Dengan begitu, perubahan pada list milik pemanggil tidak ikut mengubah isi keranjang, dan sebaliknya. Ini prinsip yang sama dengan getter yang mengembalikan salinan di Bab 2.

\subsection{Salinan Dangkal dan Salinan Dalam}

Method \texttt{copy()} dan fungsi \texttt{copy.copy()} membuat \emph{salinan dangkal} (\emph{shallow copy}): wadah luarnya baru, tetapi isinya masih object yang sama. Untuk struktur bersarang, seperti list berisi list, gunakan \texttt{copy.deepcopy()} yang menyalin seluruh isi hingga ke dalam.

\begin{kode}[shallow copy vs deep copy]
import copy

kelas = [["Ani", 90], ["Budi", 85]]
dangkal = copy.copy(kelas)
dalam = copy.deepcopy(kelas)

kelas[0][1] = 0          # ubah nilai Ani di data asli
print(dangkal[0])        # ikut berubah: isi masih dipakai bersama
print(dalam[0])          # aman: seluruh isi ikut disalin
\end{kode}

\begin{keluaran}
['Ani', 0]
['Ani', 90]
\end{keluaran}

\begin{penting}[Frozen Tidak Berarti Membeku Seluruhnya]
Data class dengan \texttt{frozen=True} hanya mencegah field \emph{diganti}. Jika sebuah field berisi list, isi list itu tetap bisa diubah dengan \texttt{append()}. Untuk data yang benar-benar tidak boleh berubah, gunakan tipe immutable di dalamnya, misalnya \texttt{tuple} sebagai pengganti \texttt{list}.
\end{penting}

% =====================================================================
\section{Duck Typing}
% =====================================================================

\begin{quote}
\itshape ``Jika ia berjalan seperti bebek dan bersuara seperti bebek, maka ia adalah bebek.''
\end{quote}

\emph{Duck typing} adalah gaya Python dalam menilai object: yang penting bukan \emph{class apa} object itu, melainkan \emph{apa yang bisa dilakukannya}. Selama object memiliki method atau attribute yang dibutuhkan, ia boleh dipakai, walaupun tidak ada hubungan pewarisan sama sekali.

\begin{kode}[Tiga class tanpa hubungan pewarisan]
class Bebek:
    def suara(self):
        return "Kwek kwek"

class Robot:
    def suara(self):
        return "Bip bup"

class Radio:
    def suara(self):
        return "Siaran berita pukul tujuh"

def bunyikan(benda):
    print(benda.suara())        # tidak peduli class-nya

for x in [Bebek(), Robot(), Radio()]:
    bunyikan(x)
\end{kode}

\begin{keluaran}
Kwek kwek
Bip bup
Siaran berita pukul tujuh
\end{keluaran}

Fungsi \texttt{bunyikan()} tidak memeriksa tipe dan tidak mensyaratkan parent tertentu. Ia hanya ``mempercayai'' bahwa benda yang diberikan bisa bersuara.

\subsection{Duck Typing di Pustaka Bawaan}

Banyak fitur Python bekerja berdasarkan duck typing. Fungsi \texttt{len()} bekerja pada object apa pun yang memiliki method \texttt{\_\_len\_\_}, dan perulangan \texttt{for} bekerja pada object apa pun yang memiliki \texttt{\_\_iter\_\_}. Dengan menambahkan dua method tersebut, class buatan kita langsung bisa dipakai seperti list.

\begin{kode}[Class sendiri yang berperilaku seperti koleksi]
class Playlist:
    def __init__(self, nama, lagu):
        self.nama = nama
        self._lagu = list(lagu)

    def __len__(self):
        return len(self._lagu)

    def __iter__(self):
        return iter(self._lagu)

pagi = Playlist("Pagi", ["Lagu A", "Lagu B", "Lagu C"])
print(len(pagi))
for judul in pagi:
    print("-", judul)
print("Lagu B" in pagi)
\end{kode}

\begin{keluaran}
3
- Lagu A
- Lagu B
- Lagu C
True
\end{keluaran}

\subsection{EAFP vs LBYL}

Ada dua gaya menghadapi object yang mungkin tidak punya kemampuan yang dibutuhkan:

\begin{itemize}
  \item \textbf{LBYL} (\emph{Look Before You Leap}): periksa dulu, baru pakai, misalnya dengan \texttt{hasattr()}.
  \item \textbf{EAFP} (\emph{Easier to Ask Forgiveness than Permission}): langsung pakai, lalu tangani kesalahannya jika gagal. Gaya ini dianggap lebih \emph{Pythonic}.
\end{itemize}

\begin{kode}[Dua gaya]
class Batu:
    pass

def bunyikan_lbyl(benda):
    if hasattr(benda, "suara"):
        print(benda.suara())
    else:
        print("(diam)")

def bunyikan_eafp(benda):
    try:
        print(benda.suara())
    except AttributeError:
        print("(diam)")

bunyikan_lbyl(Batu())
bunyikan_eafp(Batu())
\end{kode}

\begin{keluaran}
(diam)
(diam)
\end{keluaran}

Blok \texttt{try}--\texttt{except} dibahas lengkap di Bab 8. Untuk saat ini, cukup pahami bahwa duck typing memberi keleluasaan besar, dengan satu kelemahan: kesalahan baru ketahuan saat program \emph{berjalan}. Subbab berikutnya membahas cara mendeteksinya lebih awal.

% =====================================================================
\section{Structural Typing}
% =====================================================================

Ada dua cara sebuah bahasa menilai kecocokan tipe:

\begin{itemize}
  \item \textbf{Nominal typing}: object cocok bila class-nya \emph{dinyatakan} sebagai turunan tipe yang diminta. Contohnya \texttt{Persegi} cocok sebagai \texttt{BangunDatar} karena ditulis \texttt{class Persegi(BangunDatar)}. Abstract class dari subbab pertama bekerja dengan cara ini.
  \item \textbf{Structural typing}: object cocok bila \emph{bentuknya} sesuai, yaitu memiliki method dan attribute yang diminta, tanpa perlu menyatakan hubungan pewarisan.
\end{itemize}

Structural typing adalah versi duck typing yang dapat diperiksa oleh alat bantu sebelum program dijalankan. Untuk itu kita perlu mengenal \emph{type hint} terlebih dahulu.

\subsection{Type Hint Dasar}

\emph{Type hint} adalah anotasi yang menyatakan tipe yang diharapkan untuk parameter, nilai kembalian, dan variabel. Kita sudah memakainya pada field data class.

\begin{kode}[Type hint]
def luas_persegi_panjang(panjang: float, lebar: float) -> float:
    return panjang * lebar

def cari_mahasiswa(nim: str, data: dict[str, str]) -> str | None:
    return data.get(nim)

nilai: list[int] = [90, 85, 78]

print(luas_persegi_panjang(4, 2.5))
print(cari_mahasiswa("002", {"001": "Ani", "002": "Budi"}))
print(luas_persegi_panjang("ab", 3))   # tetap jalan: hint tidak dipaksakan
\end{kode}

\begin{keluaran}
10.0
Budi
ababab
\end{keluaran}

Baris terakhir menunjukkan hal penting: \textbf{Python tidak memeriksa type hint saat program berjalan}. Type hint adalah dokumentasi yang dapat dibaca manusia, editor (untuk saran otomatis dan peringatan), serta alat pemeriksa statis seperti \texttt{mypy}.

\begin{table}[htbp]
\caption{Penulisan type hint yang sering dipakai.}
\begin{tabularx}{\linewidth}{P{44mm}L}
\kepala \kp{Type hint} & \kp{Artinya} \\
\texttt{int}, \texttt{str}, \texttt{float}, \texttt{bool} & tipe dasar \\
\texttt{list[int]} & list berisi bilangan bulat \\
\texttt{dict[str, float]} & dictionary dengan kunci string dan nilai float \\
\texttt{tuple[int, str]} & tuple dua elemen: int lalu str \\
\texttt{str | None} & string atau \texttt{None} (Python 3.10 ke atas) \\
\texttt{-> None} & fungsi tidak mengembalikan nilai \\
\texttt{Mahasiswa} & object dari class \texttt{Mahasiswa} atau turunannya \\
\end{tabularx}
\end{table}

\subsection{\texttt{typing.Protocol}}

\texttt{Protocol} dari modul \texttt{typing} memungkinkan kita mendefinisikan sebuah ``bentuk'' yang diharapkan. Class lain dianggap memenuhi protocol tersebut \textbf{tanpa perlu mewarisinya}, selama memiliki method dengan nama dan bentuk yang sama.

\begin{kode}[suara.py]
from typing import Protocol

class BisaBersuara(Protocol):
    def suara(self) -> str: ...

class Bebek:                       # tidak mewarisi BisaBersuara
    def suara(self) -> str:
        return "Kwek kwek"

class Robot:
    def suara(self) -> str:
        return "Bip bup"

class Batu:
    pass

def bunyikan(benda: BisaBersuara) -> None:
    print(benda.suara())

bunyikan(Bebek())
bunyikan(Robot())
bunyikan(Batu())                   # akan ditandai oleh mypy
\end{kode}

Saat dijalankan biasa, dua baris pertama berhasil dan baris ketiga gagal dengan \texttt{AttributeError}. Namun dengan pemeriksa statis \texttt{mypy}, kesalahan itu sudah terdeteksi \emph{sebelum} program dijalankan.

\begin{keluaran}[Hasil pip install mypy, lalu mypy suara.py]
suara.py:22: error: Argument 1 to "bunyikan" has incompatible type "Batu"; expected "BisaBersuara"  [arg-type]
suara.py:22: note: "Batu" is missing following "BisaBersuara" protocol member:
suara.py:22: note:     suara
Found 1 error in 1 file (checked 1 source file)
\end{keluaran}

\begin{info}[Tentang Keluaran mypy]
\texttt{mypy} adalah alat terpisah yang dipasang dengan \texttt{pip install mypy}. Susunan pesannya dapat sedikit berbeda antarversi, tetapi intinya sama: nama class, protocol yang diharapkan, dan member yang belum ada. Editor seperti VS Code (dengan ekstensi Python/Pylance) menampilkan peringatan serupa langsung di layar sebagai garis bawah berwarna.
\end{info}

\subsection{Protocol yang Dapat Diperiksa Saat Program Berjalan}

Secara bawaan, protocol hanya berguna bagi pemeriksa statis. Jika ingin memakainya dengan \texttt{isinstance()}, tambahkan decorator \texttt{@runtime\_checkable}. Pemeriksaan ini hanya melihat \emph{keberadaan} method, bukan parameter maupun tipe kembaliannya.

\begin{kode}[runtime\_checkable]
from typing import Protocol, runtime_checkable

@runtime_checkable
class BisaBersuara(Protocol):
    def suara(self) -> str: ...

class Bebek:
    def suara(self) -> str:
        return "Kwek kwek"

class Batu:
    pass

print(isinstance(Bebek(), BisaBersuara))
print(isinstance(Batu(), BisaBersuara))
\end{kode}

\begin{keluaran}
True
False
\end{keluaran}

\begin{figure}[htbp]
\centering
\begin{tikzpicture}[font=\sffamily\small,
    tipe/.style={rectangle,draw=marun,line width=1pt,fill=merahmuda!50,align=center,minimum width=34mm,minimum height=10mm},
    kls/.style={rectangle,fill=marun,text=white,align=center,minimum width=26mm,minimum height=10mm}]
  \node[font=\sffamily\small\bfseries,text=merah] at (0,2.9) {Nominal (ABC)};
  \node[tipe,dashed] (abc) at (0,1.8) {\textit{BangunDatar}\\{\ttfamily\footnotesize luas(), keliling()}};
  \node[kls] (p) at (0,-0.4) {Persegi};
  \draw[turunan,marun] (p.north) -- (abc.south);
  \node[font=\sffamily\footnotesize,text=abu,align=center] at (0,-1.5) {cocok karena \textbf{dinyatakan}\\mewarisi \texttt{BangunDatar}};
  \draw[garis,line width=0.6pt] (3.6,-2.0) -- (3.6,3.2);
  \node[font=\sffamily\small\bfseries,text=merah] at (7.6,2.9) {Struktural (Protocol)};
  \node[tipe] (pr) at (7.6,1.8) {BisaBersuara\\{\ttfamily\footnotesize suara() -> str}};
  \node[kls] (be) at (5.7,-0.4) {Bebek\\{\ttfamily\footnotesize suara()}};
  \node[kls,fill=jingga] (ba) at (9.5,-0.4) {Batu\\{\ttfamily\footnotesize -}};
  \draw[-{Latex},line width=1pt,hijau,dashed] (be.north) -- node[left,font=\sffamily\footnotesize]{bentuk cocok} (pr.south west);
  \draw[-{Latex},line width=1pt,jingga,dashed] (ba.north) -- node[right,font=\sffamily\footnotesize]{tidak cocok} (pr.south east);
  \node[font=\sffamily\footnotesize,text=abu,align=center] at (7.6,-1.5) {tidak ada pewarisan;\\yang dilihat hanya \textbf{bentuknya}};
\end{tikzpicture}
\caption{Nominal typing menilai dari pernyataan pewarisan; structural typing menilai dari bentuk object.}
\end{figure}

\begin{table}[htbp]
\caption{Abstract class dibandingkan Protocol.}
\begin{tabularx}{\linewidth}{P{30mm}LL}
\kepala \kp{Aspek} & \kp{Abstract class (\texttt{ABC})} & \kp{\texttt{Protocol}} \\
Jenis typing & nominal & struktural \\
Wajib mewarisi? & ya & tidak \\
Kapan dicek & saat object dibuat (program berjalan) & oleh pemeriksa statis sebelum program berjalan \\
Boleh berisi method konkret & ya, dan diwarisi child & boleh, tetapi tidak diwarisi class yang hanya ``cocok bentuk'' \\
Cocok untuk & hierarki class milik sendiri yang berbagi kode & menerima object dari mana saja, termasuk class pustaka lain yang tidak bisa diubah \\
\end{tabularx}
\end{table}

% =====================================================================
\section{Rangkuman}
% =====================================================================

\begin{itemize}
  \item \textbf{Abstract class} (mewarisi \texttt{ABC}) berisi \textbf{abstract method} (\texttt{@abstractmethod}) yang wajib diimplementasikan child. Abstract class dan child yang kontraknya belum lengkap tidak bisa dibuat object-nya.
  \item \textbf{Data class} (\texttt{@dataclass}) membangkitkan \texttt{\_\_init\_\_}, \texttt{\_\_repr\_\_}, dan \texttt{\_\_eq\_\_} dari daftar field. Opsi penting: \texttt{order}, \texttt{frozen}, \texttt{field(default\_factory=\ldots)}, dan \texttt{\_\_post\_init\_\_}.
  \item \textbf{Mutable} berarti isi object dapat diubah. Variabel adalah label; \texttt{b = a} tidak menyalin object. Hindari default argument mutable dan bedakan salinan dangkal dari salinan dalam.
  \item \textbf{Duck typing}: object dinilai dari kemampuannya, bukan class-nya. Method seperti \texttt{\_\_len\_\_} dan \texttt{\_\_iter\_\_} membuat class sendiri bekerja dengan \texttt{len()} dan \texttt{for}.
  \item \textbf{Structural typing} dengan \texttt{typing.Protocol} memformalkan duck typing sehingga bisa diperiksa oleh \texttt{mypy}. Type hint tidak dipaksakan saat program berjalan.
\end{itemize}

% =====================================================================
\section{Kuis Teoritis}
% =====================================================================

\subsection{Pilihan Ganda}

Pilihlah satu jawaban yang paling tepat.

\begin{enumerate}
  \item Apa yang terjadi bila kita mencoba membuat object dari class yang mewarisi \texttt{ABC} dan masih memiliki abstract method?
  \begin{enumerate}[label=\alph*.,itemsep=0pt]
    \item Object berhasil dibuat, abstract method mengembalikan \texttt{None}
    \item Terjadi \texttt{TypeError}
    \item Terjadi \texttt{AttributeError} saat method dipanggil
    \item Python otomatis mengisi method tersebut
  \end{enumerate}

  \item Pernyataan yang benar tentang abstract class adalah \ldots
  \begin{enumerate}[label=\alph*.,itemsep=0pt]
    \item semua method-nya harus abstract
    \item tidak boleh memiliki constructor
    \item boleh memiliki method konkret yang memanggil abstract method
    \item hanya boleh diwarisi oleh satu child class
  \end{enumerate}

  \item Method yang \textbf{tidak} dibuat otomatis oleh \texttt{@dataclass} tanpa opsi tambahan adalah \ldots
  \begin{enumerate}[label=\alph*.,itemsep=0pt]
    \item \texttt{\_\_init\_\_}
    \item \texttt{\_\_repr\_\_}
    \item \texttt{\_\_eq\_\_}
    \item \texttt{\_\_lt\_\_}
  \end{enumerate}

  \item Cara yang benar memberi default list kosong pada field data class adalah \ldots
  \begin{enumerate}[label=\alph*.,itemsep=0pt]
    \item \texttt{nilai: list = []}
    \item \texttt{nilai: list = list}
    \item \texttt{nilai: list = field(default\_factory=list)}
    \item \texttt{nilai: list = None[]}
  \end{enumerate}

  \item Perhatikan kode berikut.
\begin{kode}[Soal 5]
x = [1, 2]
y = x
y += [3]
z = x.copy()
z.append(4)
print(x)
\end{kode}
  Keluaran kode tersebut adalah \ldots
  \begin{enumerate}[label=\alph*.,itemsep=0pt]
    \item \texttt{[1, 2]}
    \item \texttt{[1, 2, 3]}
    \item \texttt{[1, 2, 3, 4]}
    \item \texttt{[1, 2, 4]}
  \end{enumerate}

  \item Tipe berikut yang bersifat \emph{immutable} adalah \ldots
  \begin{enumerate}[label=\alph*.,itemsep=0pt]
    \item \texttt{list}
    \item \texttt{dict}
    \item \texttt{set}
    \item \texttt{tuple}
  \end{enumerate}

  \item Perhatikan kode berikut.
\begin{kode}[Soal 7]
def tambah(item, wadah=[]):
    wadah.append(item)
    return wadah

tambah("a")
print(tambah("b"))
\end{kode}
  Keluaran kode tersebut adalah \ldots
  \begin{enumerate}[label=\alph*.,itemsep=0pt]
    \item \texttt{['b']}
    \item \texttt{['a', 'b']}
    \item \texttt{['a']}
    \item \texttt{[]}
  \end{enumerate}

  \item Agar object dari class buatan sendiri dapat dipakai dengan fungsi \texttt{len()}, class tersebut perlu memiliki method \ldots
  \begin{enumerate}[label=\alph*.,itemsep=0pt]
    \item \texttt{\_\_size\_\_}
    \item \texttt{\_\_len\_\_}
    \item \texttt{\_\_count\_\_}
    \item \texttt{length()}
  \end{enumerate}

  \item Apa yang terjadi saat fungsi \texttt{def f(x: int) -> int} dipanggil dengan \texttt{f("teks")} lalu dijalankan biasa dengan \texttt{python}?
  \begin{enumerate}[label=\alph*.,itemsep=0pt]
    \item Python menolak dengan \texttt{TypeError} sebelum fungsi dijalankan
    \item Python otomatis mengubah \texttt{"teks"} menjadi bilangan
    \item Fungsi tetap dijalankan; type hint tidak dipaksakan saat runtime
    \item Program tidak bisa dijalankan sama sekali
  \end{enumerate}

  \item Perbedaan utama \texttt{Protocol} dibandingkan abstract class adalah \ldots
  \begin{enumerate}[label=\alph*.,itemsep=0pt]
    \item class yang cocok dengan protocol tidak perlu mewarisinya
    \item protocol tidak boleh memiliki method
    \item protocol hanya bisa dipakai untuk data class
    \item protocol memaksa pengecekan saat object dibuat
  \end{enumerate}
\end{enumerate}

\subsection{Uraian}

\begin{enumerate}
  \item Jelaskan mengapa memberi nilai default \texttt{return 0} pada method parent lebih berbahaya daripada menjadikannya abstract method. Berikan contoh kasus selain bangun datar.
  \item Bandingkan penulisan class pembawa data secara manual dengan \texttt{@dataclass}. Sebutkan dua situasi ketika data class \emph{kurang} tepat digunakan.
  \item Jelaskan apa yang dimaksud ``variabel adalah label, bukan kotak'' dan kaitkan dengan bug default argument mutable pada class \texttt{Keranjang}.
  \item Jelaskan perbedaan salinan dangkal dan salinan dalam. Pada kondisi apa keduanya memberikan hasil yang sama?
  \item Jelaskan hubungan antara duck typing dan structural typing. Mengapa \texttt{Protocol} berguna jika Python sendiri sudah mendukung duck typing?
\end{enumerate}

% =====================================================================
\section{Praktik Mandiri}
% =====================================================================

Kerjakan praktik berikut secara berurutan. Simpan setiap praktik dalam file \texttt{.py} terpisah, jalankan, dan pastikan keluarannya sesuai dengan contoh.

\subsection{Praktik 1: Metode Pembayaran}

Buat abstract class \texttt{Pembayaran} dengan ketentuan:

\begin{itemize}
  \item Abstract property \texttt{nama\_metode} dan abstract method \texttt{biaya\_admin(jumlah)}.
  \item Method konkret \texttt{struk(jumlah)} yang menampilkan nama metode, jumlah, biaya admin, dan total, dengan memanggil kedua anggota abstract tersebut.
  \item Tiga child: \texttt{Tunai} (tanpa biaya), \texttt{TransferBank} (biaya tetap Rp2.500), dan \texttt{EWallet} (biaya 1\% dari jumlah, minimal Rp1.000).
  \item Buat juga class \texttt{QRIS(Pembayaran)} yang sengaja hanya mengimplementasikan \texttt{nama\_metode}, lalu tunjukkan pesan kesalahan yang muncul saat membuat object-nya.
\end{itemize}

\begin{keluaran}[Contoh keluaran Praktik 1]
Tunai         Rp50.000 + admin Rp0     = Rp50.000
Transfer Bank Rp50.000 + admin Rp2.500 = Rp52.500
E-Wallet      Rp50.000 + admin Rp1.000 = Rp51.000
E-Wallet      Rp250.000 + admin Rp2.500 = Rp252.500
TypeError: Can't instantiate abstract class QRIS without an implementation for abstract method 'biaya_admin'
\end{keluaran}

\subsection{Praktik 2: Papan Peringkat dengan Data Class}

\begin{itemize}
  \item Buat data class \texttt{Peserta} dengan field \texttt{nama}, \texttt{asal}, dan \texttt{skor\_soal} (list berisi skor tiap soal, default list kosong yang aman).
  \item Tambahkan field \texttt{total} yang dihitung di \texttt{\_\_post\_init\_\_} dan tidak menjadi parameter constructor. Tolak skor di luar rentang 0--100 dengan \texttt{ValueError}.
  \item Buat data class \texttt{frozen} bernama \texttt{Lomba} berisi \texttt{nama} dan \texttt{tahun}, lalu tunjukkan bahwa field-nya tidak bisa diubah.
  \item Urutkan peserta berdasarkan \texttt{total} dari terbesar, lalu tampilkan peringkatnya. Gunakan \texttt{sorted(\ldots, key=\ldots)} atau opsi \texttt{order=True} dengan \texttt{field(compare=False)}; jelaskan pilihan Anda dalam komentar.
\end{itemize}

\begin{keluaran}[Contoh keluaran Praktik 2]
Lomba(nama='Olimpiade Pemrograman', tahun=2026)
1. Citra  (UNESA)  total 275
2. Ani    (ITS)    total 260
3. Budi   (UNAIR)  total 245
FrozenInstanceError: cannot assign to field 'tahun'
\end{keluaran}

\subsection{Praktik 3: Ekspor Laporan dengan Protocol}

\begin{langkah}{Langkah Pengerjaan}
  \item Definisikan \texttt{Protocol} bernama \texttt{BisaDiekspor} dengan method \texttt{ekspor(self) -> str} dan attribute \texttt{judul: str}.
  \item Buat tiga class yang \textbf{tidak saling mewarisi}: \texttt{LaporanNilai}, \texttt{GrafikKehadiran}, dan \texttt{TabelJadwal}. Masing-masing mengimplementasikan \texttt{ekspor()} dengan format berbeda (misalnya teks, CSV, atau baris tabel).
  \item Buat class \texttt{Arsip} yang menyimpan banyak object \texttt{BisaDiekspor}, mengimplementasikan \texttt{\_\_len\_\_} dan \texttt{\_\_iter\_\_}, serta memiliki method \texttt{ekspor\_semua()}. Pastikan constructor-nya aman dari jebakan default mutable.
  \item Tambahkan satu class \texttt{Catatan} yang tidak memiliki \texttt{ekspor()}, lalu masukkan ke arsip. Jalankan \texttt{mypy} pada file Anda dan salin pesan kesalahannya ke dalam komentar. Jika tidak bisa memasang \texttt{mypy}, gunakan \texttt{@runtime\_checkable} dan \texttt{isinstance()} untuk menolak object yang tidak sesuai.
\end{langkah}

\begin{penting}[Kriteria Penilaian]
\begin{itemize}
  \item Kontrak abstract class dan protocol dirancang dengan tepat dan benar-benar ditegakkan.
  \item Data class memakai default mutable yang aman serta validasi di \texttt{\_\_post\_init\_\_}.
  \item Tidak ada data yang tanpa sengaja dipakai bersama antar-object.
  \item Type hint ditulis pada semua parameter dan nilai kembalian di Praktik 3.
\end{itemize}
\end{penting}